%!TEX root = ../../main.tex
\section{주 결과: 소규모 교전 코호트 S}
\label{sec:pred-primary-S}

표 \ref{tab:primary-S}는 15.16 TEST의 S 행(\numSTest{} 교전, \numSTestMatches{} 경기)에서 항등 보정기 아래 $\qS$, $\PTflexS$, $\PTlinS$의 점수를 보인다. 주 대비 $\dBrier(\qS - \PTflexS)$는
\begin{equation}
  \dBrier = \dBrierS, \qquad 95\%\ \text{CI}\ \ci{\dBrierSLo}{\dBrierSHi}
\end{equation}
이며 $\dBrier > 0$인 추출은 없었다. $\PTlinS$에 대한 연속성은 \dBrierSlin{} \ci{\dBrierSlinLo}{\dBrierSlinHi}이다. 민감도로, RR12의 두 단계 보정기 선택(모든 S 모델에 양의 기울기 시그모이드)을 적용하면 주 대비는 \dBrierSsens{} \ci{\dBrierSsensLo}{\dBrierSsensHi}로 부호와 크기가 계약대로의 항등 결과와 일치한다.

%% table 객체: tab:primary-S — 출처 SCALE_SPLIT_TvsS_RESULTS json table1_S_TEST_identity, table2a_identity, table2b_sensitivity; RR12_S_qS_id json
\begin{table}[!t]
  \centering
  \caption{소규모 교전 코호트 S의 주 결과 (15.16 TEST)}
  \label{tab:primary-S}
  \small
  \begin{tabular}{@{}lrrr@{}}
    \toprule
    모델 & Brier $\downarrow$ & 로그 손실 $\downarrow$ & AUC $\uparrow$ \\
    \midrule
    $\qS$ (항등) & \qSBrier{} & \qSLogloss{} & \qSAUC{} \\
    $\PTflexS$ & \ptflexSBrier{} & \ptflexSLogloss{} & \ptflexSAUC{} \\
    $\PTlinS$ & \ptlinSBrier{} & \ptlinSLogloss{} & \ptlinSAUC{} \\
    \midrule
    \multicolumn{4}{@{}l}{\textbf{주 대비} $\dBrier(\qS - \PTflexS)$ = \textbf{\dBrierS{}}, 95\% CI \ci{\dBrierSLo}{\dBrierSHi}} \\
    \multicolumn{4}{@{}l}{민감도 (RR12 보정기 선택, 시그모이드): \dBrierSsens{}, 95\% CI \ci{\dBrierSsensLo}{\dBrierSsensHi}} \\
    \multicolumn{4}{@{}l}{연속성 $\dBrier(\qS - \PTlinS)$ = \dBrierSlin{}, 95\% CI \ci{\dBrierSlinLo}{\dBrierSlinHi}} \\
    \bottomrule
  \end{tabular}
  \tabnote{$n = \numSTest{}$ 교전, \numSTestMatches{} 경기; 경기 가중; 항등 보정기; 구간은 \numBootDraws{}회 경기 군집 백분위 부트스트랩(시드 \bootSeed).}
\end{table}


각 코호트 안에서 교전 전 상태 모델은 시험된 $(\ppre, t)$의 유연한 함수보다 비슷한 절대 폭만큼 적정 점수 손실을 낮춘다. 두 블록은 행 수(\numSTest{} 대 \numTTest), 경기 가중 양성 비율(S \posRateS, T \posRateT), $\Bforty$ 비율(\shareSBforty\% 대 \shareTBforty\%)이 다르므로 절대 Brier와 AUC는 같은 척도 위에 있지 않으며 비교하지 않는다. 본 논문은 두 교전 규모의 상대적 예측 가능성에 관해 어떤 주장도 하지 않는다.
